\clearpage
\SourceLeaf{377}%
\begin{center}\large [CONFORME ABBILDUNG. KRÜMMUNGSMASS.]\end{center}
\begin{center}[1.]\end{center}
\begin{center}\textsc{Gauss} an \textsc{Schumacher}. Göttingen, 5. Juli 1816.\end{center}

\mbox{. . . . . . .} Das Programm mit der Preisfrage Ihrer Societät ist mir noch nicht zu Gesichte gekommen. Mit \textsc{Lindenau} habe ich auch über eine Preisfrage conferirt, die in der neuen Zeitschrift mit dem Preise von 100 Ducaten aufgegeben werden soll. Mir war eine interessante Aufgabe eingefallen, nemlich:
\begin{quote}
\noindent »\textbf{allgemein} eine gegebene Fläche so auf einer andern (gegebenen) zu\\
»projiciren (abzubilden), dass das Bild dem Original in den kleinsten\\
»Theilen ähnlich werde«.
\end{quote}

Ein specieller Fall ist, wenn die erste Fläche eine Kugel, die zweite eine Ebene ist. Hier sind die stereographische und die merkatorische Projectionen particuläre Auflösungen. Man will aber die \textls[120]{allgemeine} Auflösung, worunter alle particulären begriffen sind, für jede Arten von Flächen.

Es soll darüber in dem Journal philomathique bereits von \textsc{Monge} und \textsc{Ponsot} gearbeitet sein (wie \textsc{Burckhardt} an \textsc{Lindenau} geschrieben hat), allein da ich nicht genau weiss wo, so habe ich noch nicht nachsuchen können, und weiss daher nicht, ob jener Herren Auflösungen ganz meiner Idee entsprechen und die Sache erschöpfen. Im entgegengesetzten Fall schiene mir diess einmal eine schickliche Preisfrage für eine Societät zu sein. Bei der hiesigen kommt die Reihe des Aufgebens nur alle 12 Jahre an mich. \mbox{. . . . . . .}
